\documentclass{article}
\usepackage[utf8]{inputenc}
\usepackage{a4wide}

\begin{document}

\section*{Flood Fill}

กำหนดหน้าจอ 2 มิติ ขนาดไม่เกิน (กว้าง x ยาวหรือสูง)\texttt{200\ x\ 100}จุด

(Pixel) พิกัดมุมบนซ้ายสุดเป็น\texttt{(0,\ 0)}กำหนดค่าสีของแต่ละ Pixel

เป็นเลขจำนวนเต็มตั้งแต่ 0 ขึ้นไปแต่ไม่เกิน 4 สิ่งที่ต้องการให้ทำคือ การเทสีหรือแทน ณ

จุดที่กำหนด\texttt{(x,\ y)}รวมถึงจุดที่อยู่ข้างกันทั้งสี่ทิศ (ไม่นับแนวทแยงมุม)

และสีเดียวกัน\ul{\textbf{ต่อเนื่อง}}ไปทั้งหมด ด้วยค่าสีใหม่ที่กำหนดคือสี\texttt{8}



\hfill\break



\textbf{ตัวอย่าง}



\textbf{Input:}



\texttt{8\ 8\ 3}{}{\texttt{//}}{\texttt{บอกขนาดหน้าจอที่ต้องการ\ 8\ x\ 8\ มีสามสี\ 0,\ 1,\ 2}}{}



\texttt{7\ 3}{\texttt{//}}{\texttt{บอกว่าเทสี\ 8\ ลงไปที่ตำแหน่ง\ (row,\ col)\ =\ (7,\ 3)}}\strut \\



\texttt{1\ 1\ 1\ 1\ 1\ 1\ 1\ 1}



\texttt{1\ 1\ 1\ 1\ 1\ 1\ 0\ 0}



\texttt{1\ 0\ 0\ 1\ 1\ 0\ 1\ 1}



\texttt{1\ 2\ 2\ 2\ 2\ 0\ 1\ 0}



\texttt{1\ 1\ 1\ 2\ 2\ 0\ 1\ 0}



\texttt{1\ 1\ 1\ 2\ 2\ 2\ 2\ 0}



\texttt{1\ 1\ 1\ 1\ 1\ 2\ 1\ 1}



\texttt{1\ 1\ 1}{\texttt{1}}\texttt{1\ 2\ 2\ 1}\strut \\



\textbf{Output:}



\textbf{\texttt{8\ 8\ 8\ 8\ 8\ 8\ 8\ 8}}



\textbf{\texttt{8\ 8\ 8\ 8\ 8\ 8}}\texttt{0\ 0}



\textbf{\texttt{8}}\texttt{0\ 0}\textbf{\texttt{8\ 8}}\texttt{0\ 1\ 1}



\textbf{\texttt{8}}\texttt{2\ 2\ 2\ 2\ 0\ 1\ 0}



\textbf{\texttt{8\ 8\ 8}}\texttt{2\ 2\ 0\ 1\ 0}



\textbf{\texttt{8\ 8\ 8}}\texttt{2\ 2\ 2\ 2\ 0}



\textbf{\texttt{8\ 8\ 8\ 8\ 8}}\texttt{2\ 1\ 1}



\textbf{\texttt{8\ 8\ 8}}{\textbf{\texttt{8}}{\textbf{\texttt{8}}\texttt{2\ 2\ 1}}}



\subsection*{Input}
ข้อมูลเข้าบรรทัดที่ 1

เป็นขนาดของหน้าจอ\texttt{(x,\ y)}ตามด้วยจำนวนสีที่มีในตอนเริ่มต้น\texttt{(M)}



ข้อมูลเข้าบรรทัดที่ 2 เป็นตำแหน่งที่เทสี\texttt{8}ลงไปที่พิกัด\texttt{(p,q)}\\



ข้อมูล{เข้า}บรรทัดที่ 3

ถึง\texttt{y\ +\ 2}เป็นตารางของเลขจำนวนเต็มแสดงค่าสี\texttt{0}ถึง\texttt{M\ –\ 1}



\hfill\break



กำหนดให้\texttt{0\ \textless{}=\ x},\texttt{p\ \textless{}\ 200},\texttt{0\ \textless{}=\ y},\texttt{q\ \textless{}\ 100},\texttt{0\ \textless{}=\ M\ \textless{}=\ 4}\\



\subsection*{Output}
ตารางใหม่หลังจากทำการเทสีลงไปแล้ว ดังตัวอย่างข้างต้น



\end{document}
